\section{From balanced cuts to circuit lower bounds}\label{sec:comparison}

Each monomial of a set-multilinear polynomial uses one variable from
each of $d$ disjoint rows. If every row has length $q$, its coefficients form a
tensor $T$ indexed by $[q]^d$, where $[q]=\{1,\ldots,q\}$. Write
\[
 f_T(X)=\sum_{j_1,\ldots,j_d\in[q]} T_{j_1\cdots j_d}
                    X_{1,j_1}\cdots X_{d,j_d}.
\]
Taking a tensor power gives more variables without increasing the
degree. In $T^{\otimes h}$ each row index is a word of length $h$, so
$f_{T^{\otimes h}}$ still has degree $d$ and now has $dq^h$ variables. Brand et al.\ study the
asymptotic circuit complexity of these grouped powers \cite{BrandMany}.

\begin{theorem}[Explicit tensors of low degree]\label{thm:lowdegree}
There are an absolute positive integer $a$ and absolute constants $c,C>0$
with the following property. For every sufficiently large even $d$, set
$m=\lceil\log_2 d\rceil$ and $w=\lceil a d m^{3/2}\rceil$. There is a
tensor $T_d$ with entries in $\{0,1\}$ on $d$ rows of common side $q_0$,
with $w\le\log q_0\le Cw$, such that over every field $F$ whose
characteristic has no ports,
\[
  \bsz_F^{[d]}(f_{T_d^{\otimes h}})
       \ge \exp(c h d w/\sqrt m) \qquad(h\ge1).
\]
An entry of $T_d$ is computable in time polynomial in $w$, uniformly in $d$.
\end{theorem}

\begin{corollary}[Bounds in the number of variables]\label{cor:variable-budget}
Over every field whose characteristic has no ports, there is a uniform
$\VNP$ family on at most $n$ variables of degree
$d=\Theta(\log n/(\log\log n)^{3/2})$ and degree-component border size
at least $\exp(\Omega((\log n)^2/(\log\log n)^2))$ for all sufficiently
large $n$. For each fixed sufficiently large even $d$ and all
$n\ge d q_0(d)$, grouped powers of the same $T_d$ give a uniform
$\VNP$ family of degree $d$, using at most $n$ variables and requiring
size $n^{\Omega(d/\sqrt{\log d})}$ in this model.
\end{corollary}

To prove Theorem~\ref{thm:lowdegree}, we extract a sum of low-rank
tensors from a small circuit. The cut on which a summand has low rank
can vary with the summand. We place the rows once and bound the cost
of every summand at those same sites. A second estimate gives an
explicit tensor whose price exceeds the total cost of this sum unless
the circuit is large. We use natural logarithms for prices.

\subsection{Balanced cuts in a circuit}

A bipartition of $[d]$ is \emph{balanced} if each side has at least
$d/3$ rows. The matrix of a tensor on such a cut has rows indexed by
the choices on one side and columns indexed by the choices on the
other. Its rank is the least number of products across the cut whose
sum is the tensor. Allowing the cut to vary over all nontrivial
bipartitions in such a sum gives Naslund's partition rank
\cite{Naslund}. A related lower bound for multiplicative complexity
uses products across three blocks of sizes $d-2,1,1$
\cite{BrandKaskiWang}. A product on one cut can have large matrix
rank on another, which limits the information from testing a single
flattening (see \cite{egow} for barriers to matrix-rank methods).
Here the balanced-gate argument of arithmetic depth reduction
\cite{VSBR} ensures that each product can use a balanced cut.

\begin{lemma}[Balanced cuts]\label{lem:balanced}
For $d\ge4$, the coefficients that use one variable in each of $d$
rows of the output of a size-$s$ circuit form a sum of at most $2^d$
tensors. Each summand has rank at most $2^{O(d)}(s+1)$ on its own
balanced cut, with at most one summand per cut. For degree-component
border size $s$, the coefficient tensor lies in the closure of the
set of these sums.
\end{lemma}

\begin{proof}
Set any variables outside the designated rows to zero. At each gate
keep the component of row multidegree $\mathbf1_I$ for
every $I\subseteq[d]$. Sums and differences act componentwise. At a
product, sum over decompositions of $I$ into two disjoint row sets.
Every discarded monomial has degree at least two in some row and
cannot contribute to a later $\mathbf1_I$ component. Computing these
components gives a circuit of size $2^{O(d)}(s+1)$ with a
row set at each gate, including the empty row set for scalar gates.

Expand the full-set output into parse trees. In each tree follow the
selected branch at a sum or difference and the child with more rows
at a product, breaking ties by a fixed rule. Stop at the first gate
with at most $2d/3$ rows. Its row set has at least $d/3$ elements,
since a product step decreases the row count by at most half and an
additive step preserves it.

Group the trees by their stopping gate. Any parse below that gate
fits each context in the group, since the path choice uses only the
row sets and branches above it. The other factors in the context use
the complementary rows. Summing the group thus gives one product
across its cut, with the signs of differences included in the
expansion. Collect products from gates with the same cut to obtain
the rank bound. Extracting these coefficients from the degree-$d$
component is a linear map, so the decomposition gives the asserted
containment after taking closures.
\end{proof}

\subsection{Two price estimates at one placement}

Work over an uncountable algebraically closed extension $K$ of
$\overline F$. A lower bound over $K$ descends to $F$. For each
circuit architecture the coefficient map is polynomial in the
constants, and the kernel of its pullback on coordinate rings
commutes with scalar extension. There are finitely many architectures
at a fixed format and size.

\paragraph{One placement for every cut.}
Theorems~\ref{prod:main} and~\ref{lamp:certificate} supply $d$ unit
sites $x_1,\ldots,x_d$ in a Hadamard metric cone, with uniform
Cartan barycenter at the tip $o$. The barycenter minimizes the mean
squared distance to the sites. The cone has dilations $z\mapsto tz$
for $t\ge0$. Write $\|z\|=d(o,z)$ and
\[
 \langle x,z\rangle=\tfrac12\bigl(\|x\|^2+\|z\|^2-d(x,z)^2\bigr).
\]
The cone distance law gives $\langle x,tz\rangle=t\langle x,z\rangle$
\cite{BH}. For a cut $\ell$, let $b_\ell(i)\in\{0,1\}$ indicate
the side containing row $i$. The same theorems give two points
$Z_{\ell,0},Z_{\ell,1}$ for each cut, at distance at most one, with
\begin{equation}\label{eq:score}
 \max_b\|Z_{\ell,b}\|=O(\sqrt d),\qquad
 \sum_i\langle x_i,Z_{\ell,b_\ell(i)}\rangle
                       \ge c_0 d/\sqrt m .
\end{equation}
The norm bound uses the fact that there are at most $2^d$ cuts.
The sum in Equation~\eqref{eq:score} measures the first-order decrease
in total distance to the rows when two suppliers move away from $o$
along their respective rays. We take the product with a Hilbert space
and choose homogeneous prices as in Section~\ref{sec:prices}.
Lemma~\ref{price:diagonal} bounds the price of sharing an $L$-valued
label among the row sites by $O(\sqrt m\log L)$.

\paragraph{A tensor that nearly attains the star bound.}
For any integer $dm\le w\le d^2$, Theorem~\ref{pack:witness} supplies
a single binary tensor $T$ of log side $\Theta(w)$. A local linear
map on each row takes $T$ to a tensor $\widehat U$ over $K$ whose
choice can depend on the priced placement. Each row of $\widehat U$
has the same finite alphabet $V$, with a real weight $t(v)$ attached
to each letter $v$. In a grouped power the weight of a word is the
sum of its letter weights. To price a tensor $A$ in these bases,
take the best sector of $A^{\otimes k}$ with prescribed word weight
in each row, adding the total prescribed weight to its ordinary
price. The limit of this maximum divided by $k$ is the weighted price
$\mathcal B_x(A)$ of Equation~\eqref{pack:weighted}.

Theorem~\ref{pack:witness} gives a function $\mathcal F$ on positive
distances with
\begin{align}
 \mathcal B_x(\widehat U)&\ge d\mathcal F(1)-O(d),&
 r\log\sum_{v\in V} e^{t(v)/r}&\le \mathcal F(r),\nonumber\\
 \mathcal F'(1)&=\Theta(w),&
 \sup_{1/2\le r\le2}|\mathcal F''(r)|&=O(w).
 \label{eq:pack-interface}
\end{align}
A supplier at the tip can send each index distance one to its row.
The resulting star has weighted price bounded by $d\mathcal F(1)$.
The witness comes within $O(1)$ per row of this benchmark.
Lemma~\ref{pack:value} gives the corresponding rules for comparing
supplies. Weighted prices add on grouped products and have closed
sublevels. Summing $L$ tensors costs at most their largest weighted
price plus the price of sharing an $L$-valued label. Sending an index
distance $r$ adds at most $r\log\sum_{v\in V}e^{t(v)/r}$ to the
ordinary price at the suppliers.

\begin{proof}[Proof of Theorem~\ref{thm:lowdegree}]
Let $s$ be the degree-component border size of $f_{T^{\otimes h}}$.
Lemma~\ref{lem:balanced} puts $T^{\otimes h}$ in the closure of
sums over balanced cuts with rank at most $2^{O(d)}(s+1)$ per cut.
Applying the local maps to obtain $\widehat U^{\otimes h}$ preserves
this conclusion. The maps are linear on coefficient arrays and can
only decrease cut ranks. The same conclusion holds if one starts
with a size-$s$ border circuit for $f_{\widehat U^{\otimes h}}$ over
$K$. We bound $s$ in either case.

For a term of rank $R\ge1$ on cut $\ell$, place two suppliers at
$y_b=\tau Z_{\ell,b}$ for a positive $\tau$ to be chosen. A rank
factorization represents the term by one wire carrying a label in
$[R]$, with one factor computed at each end. Since the suppliers
are at distance at most $\tau$, the wire costs at most $\tau\log R$.
Send each grouped index from its supplier to its row. Its sum of
exponential weights factors over the $h$ letters of a word.
Equation~\eqref{eq:pack-interface} therefore bounds the term's
weighted price by
\begin{equation}\label{eq:cut-price}
 \tau\log(s+1)+O(\tau d)+h\sum_{i=1}^d \mathcal F(r_i),\qquad
 r_i=d(x_i,y_{b_\ell(i)}).
\end{equation}
Figure~\ref{fig:supply} compares the two networks.
\begin{figure}[tbp]
\centering
\begin{tikzpicture}[x=1cm,y=1cm,font=\small,
  row zero/.style={circle,draw=ink,fill=ink,inner sep=0pt,minimum size=4.8pt},
  row one/.style={rectangle,draw=sea,fill=white,inner sep=0pt,minimum size=4.8pt},
  supplier/.style={circle,draw=ink,fill=white,inner sep=0pt,minimum size=7pt,line width=.9pt},
  leg/.style={draw=ink!60,line width=.55pt},
  rank wire/.style={draw=rust,line width=1.5pt}]
\foreach \panel/\shift in {star/0,cut/5.8}{
  \begin{scope}[xshift=\shift cm]
    \foreach \i/\xx/\yy in {1/1.4/.90,2/1.85/0,3/1.4/-.90}{
      \coordinate (\panel-l\i) at (-\xx,\yy);
      \coordinate (\panel-r\i) at (\xx,\yy);
    }
  \end{scope}
}
\node[text=ink] at (0,1.47) {\textbf{One supplier}};
\node[text=ink] at (5.8,1.47) {\textbf{Two suppliers for a cut}};
\coordinate (origin) at (0,0);
\coordinate (supplier0) at (5.04,0);
\coordinate (supplier1) at (6.56,0);
\foreach \i in {1,2,3}{
  \draw[leg] (origin) -- (star-l\i);
  \draw[leg] (origin) -- (star-r\i);
  \draw[leg] (supplier0) -- (cut-l\i);
  \draw[leg] (supplier1) -- (cut-r\i);
}
\draw[rank wire] (supplier0) -- (supplier1)
   node[midway,above=3pt,font=\footnotesize,text=rust] {$R$ labels};
\path (origin) -- (star-r2) node[midway,above=1pt] {$1$};
\path (supplier1) -- (cut-r2) node[midway,above=1pt] {$r_i$};
\foreach \panel in {star,cut}{
  \foreach \i in {1,2,3}{
    \node[row zero] at (\panel-l\i) {};
    \node[row one] at (\panel-r\i) {};
  }
  \node[right=2pt] at (\panel-r2) {$x_i$};
}
\node[supplier,label=below:$o$] at (origin) {};
\node[supplier,label=below:$y_0$] at (supplier0) {};
\node[supplier,label=below:$y_1$] at (supplier1) {};
\end{tikzpicture}
\caption{Two supplies on the same row sites, drawn schematically for one
copy. Filled circles and open squares indicate the two sides of a balanced cut.
At the tip $o$, the total cost of the legs is at most $d\mathcal F(1)$.
Placing a supplier $y_b$ on each side reduces the sum of the leg bounds
to $\sum_i\mathcal F(r_i)$, where $r_i$ is the distance from $x_i$ to
its supplier. The rank index travels on the wire joining the suppliers
at a cost of at most $\tau\log R$.}
\label{fig:supply}
\end{figure}


Set $\tau=\eta/(d\sqrt m)$ for a sufficiently small absolute $\eta>0$.
The cone distance law at a unit site gives
$r_i=1-\tau\langle x_i,Z_{\ell,b_\ell(i)}\rangle+O(\tau^2d)$
and $|r_i-1|=O(\tau\sqrt d)$. Using the two derivative estimates of
Equation~\eqref{eq:pack-interface} and the sum in Equation~\eqref{eq:score}
gives
\[
 \sum_i \mathcal F(r_i)\le d\mathcal F(1)
                 -c_0\mathcal F'(1)d\tau/\sqrt m+O(w\tau^2d^2)
               \le d\mathcal F(1)-c_1 w/m .
\]
The linear term saves order $\eta w/m$. Its error is $O(\eta^2w/m)$,
so the same choice of $\eta$ gives the saving on every balanced cut.

There are at most $2^d$ terms. Summing them costs $O(d\sqrt m)$ once,
even for $h$ copies, and closedness preserves the resulting bound on
$\widehat U^{\otimes h}$. Its weighted price is at least
$h(d\mathcal F(1)-O(d))$ by additivity. Cancel the common contribution
$hd\mathcal F(1)$ to obtain
\begin{equation}\label{eq:balance-cost}
    \tau\log(s+1)+O(hd+d\sqrt m)\ \ge\ c_1 h w/m .
\end{equation}
Take $w=\lceil a d m^{3/2}\rceil$. For sufficiently large absolute
$a$ the saving $c_1 h w/m$ absorbs both errors for every $h\ge1$,
and $dm\le w\le d^2$ for all sufficiently large $d$. Dividing by
$\tau$ yields
\begin{equation}\label{eq:packed-hardness}
 \bsz_K^{[d]}(f_{\widehat U^{\otimes h}})
       \ge \exp(c_2 h d w/\sqrt m)
\end{equation}
and the same bound for $f_{T^{\otimes h}}$ over $F$.
The entry algorithm and the side bounds follow from
Theorem~\ref{pack:witness}.
\end{proof}

\subsection{From the explicit tensor to complete polynomials}

\begin{proof}[Proof of Corollary~\ref{cor:variable-budget}]
The construction in Theorem~\ref{pack:witness} pads the row alphabets
of $T_d$ with zeros to the predetermined size
$q_0(d)=2^{\lceil C'w\rceil}$ for a fixed absolute integer $C'$.
For a variable budget $n$, take the largest even $d$ with
$d q_0(d)\le n$. Maximality gives $n<(d+2)q_0(d+2)$, and increasing
$d$ to $d+2$ changes $w$ by at most a constant factor. Thus
$w=\Theta(\log n)$ and
$d=\Theta(\log n/(\log\log n)^{3/2})$. At $h=1$,
Theorem~\ref{thm:lowdegree} gives
\begin{equation}\label{eq:vnphard}
   \log \bsz_F^{[d]}(f_{T_d})
        =\Omega\!\left(\frac{(\log n)^2}{(\log\log n)^2}\right).
\end{equation}
For fixed sufficiently large even $d$ and $n\ge d q_0(d)$, use the
largest $h$ such that $d q_0(d)^h\le n$. Then
$h\log q_0(d)\ge\tfrac12\log(n/d)$ and $n\ge d^2$, giving
$n^{\Omega(d/\sqrt{\log d})}$.

For uniform $\VNP$ membership, guess the address of one variable in
each row, compute the tensor entry, and accept if it is one. In a
grouped power, the entry is the product of the individual bits. Each
accepting address contributes the product of its selected variables.
The Boolean computation has size polynomial in $n$. To express it
as a Valiant sum \cite{Valiant}, also guess its gate values and check
each bounded-fan-in gate by a constant-degree truth-table polynomial.
Each address has a unique consistent assignment of gate values.
The same construction supplies one output bit per address in each
row, indicating which variable to use. The resulting summand has
polynomial degree and circuit size. These Boolean tests work over
every characteristic.
\end{proof}

The completeness projections of Valiant \cite{Valiant,Burgisser}
transfer the lower bound to the two cycle polynomials. We give the
weighted verifier construction here because Part~\ref{part:port}
also needs control of degrees and of the arcs carrying variables.

\subsubsection{Preserving the variables in completeness projections}\label{sec:gadgets}

We use an explicit form of Valiant's completeness reductions
\cite{Valiant} to control the order of the graph and the placement of
the variables. A \emph{weighted Boolean verifier} is a circuit of fixed
bounded fan-in that takes Boolean guesses, tests acceptance, and computes
a list of output bits. Each listed bit has a weight in $F$ or a variable
weight. The verifier sums, over accepted guesses, the product of the
weights of the listed bits that are true. Each variable occurs once in
the list. Let $s$ count guesses, gates, wires and listed bits.

\begin{lemma}[Sparse completeness projections]\label{gadget:verifier}
The value of a weighted Boolean verifier of length $s$ is a projection
of $\HC_M$ with $M=O(s+1)$. Its directed graph has maximum indegree
and outdegree at most three. Each variable labels a single arc, and
these arcs have disjoint endpoints. If $\operatorname{char} F\ne2$,
the value is also a projection of $\Per_M$ for an order $M=O(s+1)$.
Its bipartite graph has maximum degree three and its variable edges
form a matching. Both constructions take polynomial time.
\end{lemma}

\begin{proof}
First encode the computation by clauses requiring exactly one true
literal. For each gate, introduce one indicator for each row of its
truth table and require exactly one indicator to hold. For each bit $x$
incident to the gate, form another clause from $\neg x$ and the
indicators of the rows in which $x=1$. The clause forces the selected
row to agree with $x$. Unit clauses fix constants and require
acceptance. Give every listed output a private copy $b$ of its bit $x$
using the clause $(\neg x,b)$. Every accepted guess then has exactly
one extension to a satisfying assignment, with its original weight.

There are $O(s+1)$ literal occurrences. To ensure that every clause
has length at most three, replace a long clause on
$\ell_1,\ldots,\ell_t$ by the clauses $(a_{j-1},\ell_j,\neg a_j)$
for $1\le j\le t$. Unit clauses fix $a_0=0$ and $a_t=1$. On Boolean
values these constraints say $a_j=a_{j-1}+\ell_j$ over the integers,
so the running sums give a unique extension precisely when one literal
is true. Finally add a private bit pinned true to cover the empty
verifier.

\smallskip\noindent\emph{Hamiltonian cycles.}
Make one bidirectional path per bit, and arrange the paths in cyclic
order. Put a connector between each path and the next, with incoming
arcs from either endpoint of the first path and outgoing arcs to either
endpoint of the next. Forward traversal represents true. On each path,
reserve one edge (a gap) per literal occurrence. The ordinary path
segments on either side of a gap have length at least two, as does a
path with no gaps. For a literal of clause $C$, write $u,v$ for the gap
endpoints in the direction that makes the literal true. Keep the arc
$v\to u$ and replace $u\to v$ by $u\to C\to v$, where all literals
in $C$ share the clause vertex $C$. Give a private ordinary forward arc
on each listed bit's path its output weight. Every other permitted arc
has weight one and all missing arcs have weight zero.

An interior vertex of an ordinary segment must use its two distinct
neighbors in a Hamiltonian cycle, since a return to one neighbor would
close a two-cycle. Thus each whole segment is traversed in one
direction. Across a gap, $v\to u$ is both the only possible incoming
arc at $u$ and the only possible outgoing arc at $v$. If either end
needs this arc, it fills both positions. Otherwise $u$ needs an
outgoing arc across the gap and $v$ needs an incoming arc there, forcing
$u\to C\to v$. In particular a cycle cannot enter a clause at one
gap and return at another. Each path therefore has a consistent
orientation, and each clause vertex is visited once precisely when
one of its literals is true. A satisfying assignment determines one
directed cycle through the paths in their cyclic order, counted once
in $\HC_M$, with the prescribed weight. Clause lengths give the degree
bound, and weighted arcs on distinct paths have disjoint endpoints.

\smallskip\noindent\emph{Perfect matchings.}
For each bit make a simple even bipartite cycle of length at least four.
Its two alternating perfect matchings represent true and false. Reserve
one distinct edge for each literal occurrence in the matching that makes
the literal true. For a listed bit, weight one further edge of the true
matching by its output weight. For each clause, create a row and a
column, with one prospective edge between them per literal. We must
couple each prospective edge to its occurrence edge so that they are
selected together.

Delete each such pair of edges and add three private rows and three
private columns with weight matrix
\[
  A=\begin{pmatrix}
     \tfrac12&\tfrac12&-\tfrac12\\
     -1&-1&1\\
     1&1&1
  \end{pmatrix}.
\]
For the first deleted edge, join its row to private column $1$ and
private row $1$ to its column. For the second edge, use private column
and row $2$ in the same way. All four connecting edges have weight one.
Fix the connecting edges used by a matching. They must occupy equally
many private rows and columns for an internal completion to exist. If
they occupy none, the total internal weight is $\operatorname{per}(A)=1$. With
one row and one column occupied, deleting them leaves a permanent of
zero for all four choices from indices $1,2$. With both rows and both
columns occupied, the remaining entry is $A_{33}=1$. Thus the only
states that contribute are those using neither deleted edge or both,
each with factor one. This cancellation uses the inverse of $2$.

After summing over the private vertices, each clause chooses one literal
and each bit chooses an alternating matching. Their choices agree
exactly for satisfying assignments. The unweighted cycle edges carry
weight one and all missing edges carry weight zero. With these entries
the permanent equals the verifier sum.

The only vertices that can have degree four are private rows and columns.
Split any such row into two rows, give each two of its old edges, and
add one column adjacent to both with unit weights. Exactly one row
matches externally, and its choice forces the match to the new column.
Use the transposed operation for columns of degree four. This preserves
the permanent and leaves the variable edges on distinct bit cycles.
Both constructions use only a constant number of vertices per bit and
per literal occurrence, hence have order $O(s+1)$.
\end{proof}


\begin{proof}[Proof of Theorem~\ref{thm:main} in the absence of ports]
Apply Lemma~\ref{gadget:verifier} to the Boolean verifier for $f_{T_d}$
above. Its size is polynomial in $n$ with absolute constants. Thus
$f_{T_d}$ is a projection of $\HC_M$ with $M\le A n^b$ for absolute
$A,b$, and also of $\Per_M$ under the same bound outside
characteristic two. To obtain any larger order, split a vertex of
the Hamiltonian-cycle instance into a forced path, or add an identity
block to the permanent instance.

A projection that substitutes constants can mix homogeneous degrees.
Suppose the complete polynomial of order $N$ lies in the closure
of degree-$N$ components of size-$s$ circuit outputs. For each such
circuit compute all homogeneous components through degree $N$ at each
gate, using addition or convolution. The isolated
degree-$N$ output is a full homogeneous approximant of size
$O((N+1)^2(s+1))$. Project that output, then select its degree-$d$
component. Both maps on coefficients are linear, so a degree-component
border circuit for the complete polynomial gives the same border
upper bound $O((N+1)^2(s+1))$ for $f_{T_d}$.

For each sufficiently large $N$, choose $n=\lfloor(N/A)^{1/b}\rfloor$
and pad the projected instance to order $N$. Equation~\eqref{eq:vnphard}
now bounds $O((N+1)^2(s+1))$ from below. Since $\log n\sim b^{-1}\log N$,
its logarithmic lower bound dominates the $O(\log N)$ cost of
isolating the degree. This gives the claimed exponent with an
absolute constant.
\end{proof}

\begin{corollary}[Sparse supports]\label{cor:sparse-support}
The lower bound of Theorem~\ref{thm:main} holds on directed graphs of
order $N$ with indegree and outdegree at most three, and, for the
permanent outside characteristic two, on bipartite graphs of degree
at most three with $N$ vertices on each side. More precisely, in each
case there exists such a graph at every sufficiently large $N$, and
each allowed arc or edge carries its own independent variable. In a
characteristic admitting ports these graphs also attain the
stretched-exponential bound of Theorem~\ref{thm:stretch}.
\end{corollary}

\begin{proof}
In either branch apply Lemma~\ref{gadget:verifier} to the verifier
used in the transfer, here for $f_{T_d}$ and in Section~\ref{sec:border}
for the signed verifier. Keep independent variables on all permitted
arcs or edges and set the other entries of the complete matrix to zero.
The resulting degree-$N$ graph polynomial still projects to the verifier
polynomial by specializing the permitted entries. Padding a directed
vertex into a forced path puts its incoming arcs at the first vertex
and its outgoing arcs at the last, preserving both degree bounds.
Padding a bipartite graph by disjoint edges preserves its maximum
degree as well. The same homogeneous extraction before specialization
therefore transfers the verifier lower bound to these graph polynomials
at every sufficiently large order.
\end{proof}
